\label{1.8}%
Considérons trois nombres $a, b, c$ dont la somme vaut $60$. Sachant que $c$
est la moyenne des deux autres et que $c - a = 1$, que vaut le triplet
$(a, b, c)$~?
\begin{enumerate}
  \item $(19,20,21)$
  \item $(19,21,20)$
  \item $(20,20,21)$
  \item $(21,20,19)$
  \item $(21,19,20)$
\end{enumerate}
%
%
\begin{proof}
%
\begin{equation}
  \begin{cases}
    \frac{a+b}{2} &= c \\
    c - a &= 1 \\
    a + b + c &= 60
  \end{cases}
\end{equation}
%
Combinons les première et dernière lignes pour faire apparaître $a+b$, comme suit:
%
\begin{equation*}
  a + b + c = a + b + \frac{a+b}{2} =  \frac{3}{2}(a+b) = 60 \quad \then \quad \boxed{a+b = 40} \,.
\end{equation*}
%
Nous pouvons maintenant obtenir $c$:
%
\begin{equation*}
  \boxed{c} = a + b + c - (a + b) = 60 - 40 = \boxed{20} \,.
\end{equation*}
%
Cela permet de calculer $a$. En effet, 
%
\begin{equation*}
  \boxed{a} = c - (c-a) = 20 - 1 = \boxed{19} \,. 
\end{equation*}
%
Finalement,
%
\begin{equation*}
  \boxed{b} = (a + b + c) - a - c  = 60 - 19 - 20 = \boxed{21} \,.
\end{equation*}
%
Réponse \textbf{B}.
\end{proof}